\subsection{折叠 \texorpdfstring{$\pi$}{pi} 通道的有限窗奇层消元、自校准抽取与 Koenigs 正规形}\label{subsec:conclusion-foldgauge-pi-window-accelerator-zeta-koenigs}
在上一小节已经把折叠标量通道上的黄金线性精度律、偶值 \(\zeta\) 首差定位与外置地址预算壁垒固定之后，定理 \ref{thm:fold-bin-gauge-constant-stirling-bernoulli-hierarchy}、定理 \ref{thm:xi-time-part60ab-gauge-constant-richardson-neville-accelerator}、结论 \ref{con:xi-time-part9l-chi2-generates-power-divergence-and-cumulants}、定理 \ref{thm:xi-time-part64a-fold-gauge-quintic-renormalization-normal-form} 与定理 \ref{thm:xi-time-part69-fold-gauge-formal-koenigs-linearization} 还可进一步压缩为如下几条终局性后果：H.111 中的奇次 Bernoulli 谱不仅允许显式有限窗加速，而且该加速在固定窗仿射对数类中是唯一且最优的，其整条算子塔具有统一有界的条件数；同一套系数又可由单一 \(\chi^2\) 极限常数完全自校准，并反过来给出偶值 \(\zeta\)-塔的有限窗抽取公式；与此同时，fold-gauge 误差的局部动力学属于非零乘子的 Koenigs 型，而与 AGM 的超吸引 B\"ottcher 型正规形严格断裂。

\begin{theorem}[折叠 \texorpdfstring{$\pi$}{pi} 估计的内生奇层消元层级]\label{thm:conclusion-foldgauge-pi-oddlayer-annihilation-hierarchy}
记
\[
L_m:=\log C_m-\log(2\pi),
\qquad
q:=\frac{\varphi}{2}\in(0,1),
\]
并令前移算子 \(T\) 作用于序列 \(f=(f_m)_{m\ge 1}\) 为
\[
(Tf)_m:=f_{m+1}.
\]
对每个整数 \(R\ge 1\)，定义
\[
P_R(z):=\prod_{r=1}^{R}\frac{z-q^{2r-1}}{1-q^{2r-1}}
=
\sum_{j=0}^{R}\beta_j^{(R)}z^j,
\qquad
\mathcal A_R:=P_R(T),
\]
并把 \(\mathcal A_R\) 作用到序列 \((\log C_m)_{m\ge 1}\) 上，即
\[
(\mathcal A_R\log C)_m
:=
\sum_{j=0}^{R}\beta_j^{(R)}\log C_{m+j}.
\]
再记
\[
a_s:=
\frac{B_{2s}}{s(2s-1)}
\bigl(\varphi^{-1}+\varphi^{2s-3}\bigr)
\qquad (s\ge 1).
\]
则存在显式非零常数
\[
K_R:=a_{R+1}P_R(q^{2R+1})
=
\frac{B_{2R+2}}{(R+1)(2R+1)}
\bigl(\varphi^{-1}+\varphi^{2R-1}\bigr)
P_R(q^{2R+1})
\neq 0
\]
使得
\[
(\mathcal A_R\log C)_m
=
\log(2\pi)
+
K_R\,q^{(2R+1)m}
+
O\!\bigl(q^{(2R+3)m}\bigr).
\]
从而由有限窗 \((C_m,\dots,C_{m+R})\) 定义的内生 \(\pi\)-估计
\[
\Pi_m^{\langle R\rangle}
:=
\frac12\exp\!\bigl((\mathcal A_R\log C)_m\bigr)
\]
满足
\[
\Pi_m^{\langle R\rangle}-\pi
=
\pi K_R\,q^{(2R+1)m}
+
O\!\bigl(q^{(2R+3)m}\bigr).
\]
亦即，纯用项目内部的连续 \(R+1\) 尺度数据 \((C_m,\dots,C_{m+R})\)，主误差指数便可由 \(1\) 提升到 \(2R+1\)。
\end{theorem}
\begin{proof}
由定理 \ref{thm:fold-bin-gauge-constant-stirling-bernoulli-hierarchy}，对任意固定 \(R\ge 1\) 都有
\[
L_m
=
\sum_{s=1}^{R+1}a_s q^{(2s-1)m}
+
O\!\bigl(q^{(2R+3)m}\bigr).
\]
由于 \(T\) 在模态 \(q^{(2s-1)m}\) 上的特征值为 \(q^{2s-1}\)，故
\[
(\mathcal A_R L)_m
=
\sum_{s=1}^{R+1}a_s P_R(q^{2s-1})q^{(2s-1)m}
+
O\!\bigl(q^{(2R+3)m}\bigr).
\]
又因为
\[
P_R(q^{2r-1})=0\qquad(1\le r\le R),
\qquad
P_R(1)=1,
\]
前 \(R\) 个奇层被严格消去，而常数项 \(\log(2\pi)\) 被保留，于是
\[
(\mathcal A_R\log C)_m
=
\log(2\pi)
+
a_{R+1}P_R(q^{2R+1})q^{(2R+1)m}
+
O\!\bigl(q^{(2R+3)m}\bigr).
\]
这就得到第一式。又因 \(q^{2R+1}\) 不等于 \(q,q^3,\dots,q^{2R-1}\)，故 \(P_R(q^{2R+1})\neq 0\)；再由偶 Bernoulli 数 \(B_{2R+2}\neq 0\)，可知 \(K_R\neq 0\)。

最后，设
\[
\varepsilon_m:=K_R q^{(2R+1)m}+O\!\bigl(q^{(2R+3)m}\bigr).
\]
则
\[
\Pi_m^{\langle R\rangle}
=
\frac12\exp\!\bigl(\log(2\pi)+\varepsilon_m\bigr)
=
\pi e^{\varepsilon_m}
=
\pi\bigl(1+\varepsilon_m+O(\varepsilon_m^2)\bigr).
\]
而 \(\varepsilon_m^2=O(q^{(4R+2)m})=O(q^{(2R+3)m})\) 对一切 \(R\ge 1\) 都成立，遂得
\[
\Pi_m^{\langle R\rangle}-\pi
=
\pi K_R q^{(2R+1)m}+O\!\bigl(q^{(2R+3)m}\bigr).
\]
证毕。
\end{proof}

\begin{corollary}[一级有限窗加速器的三奇阶提升]\label{cor:conclusion-foldgauge-pi-first-window-cubic-lift}
\leanverified{paper\_conclusion\_foldgauge\_pi\_first\_window\_cubic\_lift}
取 \(R=1\)。则
\[
P_1(z)=\frac{z-q}{1-q},
\qquad
(\mathcal A_1\log C)_m
=
\frac{\log C_{m+1}-q\log C_m}{1-q},
\]
从而
\[
\Pi_m^{\langle1\rangle}
:=
\frac12\exp\!\left(\frac{\log C_{m+1}-q\log C_m}{1-q}\right).
\]
并且有精确误差式
\[
\Pi_m^{\langle1\rangle}-\pi
=
\pi\,\frac{3+\sqrt5}{288}\,q^{3m}
+
O(q^{5m}).
\]
因此单次一级奇层消元已经把误差阶由 \(q^m\) 提升到 \(q^{3m}\)。
\end{corollary}
\begin{proof}
由定理 \ref{thm:conclusion-foldgauge-pi-oddlayer-annihilation-hierarchy}，
\[
K_1
=
a_2\,\frac{q^3-q}{1-q}.
\]
又由
\[
a_2=
\frac{B_4}{2\cdot 3}(\varphi^{-1}+\varphi)
=
-\frac{\sqrt5}{180},
\qquad
\frac{q^3-q}{1-q}=-q(q+1),
\]
可得
\[
K_1
=
\frac{q(q+1)\sqrt5}{180}
=
\frac{3+\sqrt5}{288}.
\]
代回定理 \ref{thm:conclusion-foldgauge-pi-oddlayer-annihilation-hierarchy} 即得所示展开。证毕。
\end{proof}

\begin{theorem}[固定窗仿射对数加速器的唯一最优性]\label{thm:conclusion-foldgauge-pi-affine-log-accelerator-uniqueness}
\leanverified{paper\_conclusion\_foldgauge\_pi\_affine\_log\_accelerator\_uniqueness}
对每个整数 \(R\ge 1\)，定义固定窗仿射对数类
\[
\mathfrak A_R
:=
\left\{
\Lambda_m=\sum_{j=0}^{R}c_j\log C_{m+j}:\ \sum_{j=0}^{R}c_j=1
\right\}.
\]
则满足
\[
\Lambda_m-\log(2\pi)=O\!\bigl(q^{(2R+1)m}\bigr)
\]
的元在 \(\mathfrak A_R\) 中唯一存在，且恰为
\[
\Lambda_m=(\mathcal A_R\log C)_m.
\]
进一步，任取 \(M<R\)。若只允许使用连续窗口
\[
(C_m,C_{m+1},\dots,C_{m+M})
\]
构成仿射对数组合，则不可能达到
\[
O\!\bigl(q^{(2R+1)m}\bigr)
\]
这一误差阶。
\end{theorem}
\begin{proof}
任取 \(\Lambda_m\in\mathfrak A_R\)，并定义
\[
Q(z):=\sum_{j=0}^{R}c_j z^j.
\]
由定理 \ref{thm:fold-bin-gauge-constant-stirling-bernoulli-hierarchy}，
\[
\Lambda_m-\log(2\pi)
=
\sum_{s\ge 1}a_s Q(q^{2s-1})q^{(2s-1)m}.
\]
若 \(\Lambda_m-\log(2\pi)=O(q^{(2R+1)m})\)，则对 \(1\le s\le R\) 必有
\[
Q(q^{2s-1})=0,
\]
否则对应的第 \(s\) 层奇次模态不会消失。另一方面，
\[
Q(1)=\sum_{j=0}^{R}c_j=1.
\]
因此 \(Q\) 是一个次数至多 \(R\) 的多项式，在 \(R\) 个互异点
\[
q,q^3,\dots,q^{2R-1}
\]
上为零，并在 \(1\) 处取值 \(1\)。由插值唯一性，
\[
Q(z)=P_R(z).
\]
于是唯一最优元正是
\[
\Lambda_m=(P_R(T)\log C)_m=(\mathcal A_R\log C)_m.
\]

若只允许使用 \(M+1\) 个连续尺度且 \(M<R\)，则相应多项式次数至多为 \(M\)。这样的多项式不可能同时在 \(R\) 个不同点
\[
q,q^3,\dots,q^{2R-1}
\]
上取零并保持 \(Q(1)=1\)。故短窗不可能达到 \(O(q^{(2R+1)m})\) 的误差阶。证毕。
\end{proof}

\begin{theorem}[加速层级的统一有界条件数]\label{thm:conclusion-foldgauge-pi-accelerator-uniform-conditioning}
把定理 \ref{thm:conclusion-foldgauge-pi-oddlayer-annihilation-hierarchy} 中的系数记作
\[
P_R(z)=\sum_{j=0}^{R}\beta_j^{(R)}z^j.
\]
则其 \(\ell^1\)-范数满足精确闭式
\[
\sum_{j=0}^{R}\bigl|\beta_j^{(R)}\bigr|
=
\prod_{r=1}^{R}\frac{1+q^{2r-1}}{1-q^{2r-1}}.
\]
特别地，无穷乘积
\[
\mathcal K_\varphi
:=
\prod_{r\ge 1}\frac{1+q^{2r-1}}{1-q^{2r-1}}
\]
存在且有限，并且
\[
\mathcal K_\varphi\approx 238.5415818448505.
\]
因此整条有限窗加速塔在仿射对数意义下具有统一有界条件数：对任意扰动序列 \(\delta=(\delta_m)_{m\ge 1}\)，都有
\[
\bigl|(\mathcal A_R\delta)_m\bigr|
\le
\left(\prod_{r=1}^{R}\frac{1+q^{2r-1}}{1-q^{2r-1}}\right)
\max_{0\le j\le R}|\delta_{m+j}|
\le
\mathcal K_\varphi \max_{0\le j\le R}|\delta_{m+j}|.
\]
\end{theorem}
\begin{proof}
记
\[
a_r:=q^{2r-1}\in(0,1).
\]
则
\[
P_R(z)
=
\prod_{r=1}^{R}\frac{z-a_r}{1-a_r}
=
\frac{1}{\prod_{r=1}^{R}(1-a_r)}
\sum_{j=0}^{R}(-1)^{R-j}e_{R-j}(a_1,\dots,a_R)z^j,
\]
其中 \(e_k\) 为基本对称多项式。由于所有 \(a_r>0\)，各 \(e_k(a_1,\dots,a_R)\) 均为正，因此系数符号严格交替，于是
\[
\sum_{j=0}^{R}\bigl|\beta_j^{(R)}\bigr|
=
\frac{\sum_{j=0}^{R}e_j(a_1,\dots,a_R)}{\prod_{r=1}^{R}(1-a_r)}
=
\frac{\prod_{r=1}^{R}(1+a_r)}{\prod_{r=1}^{R}(1-a_r)}
=
\prod_{r=1}^{R}\frac{1+a_r}{1-a_r}.
\]
这就得到有限阶闭式。

再由
\[
\sum_{r\ge 1}a_r=\sum_{r\ge 1}q^{2r-1}=\frac{q}{1-q^2}<\infty
\]
以及展开
\[
\log\frac{1+a_r}{1-a_r}=2a_r+O(a_r^3),
\]
可知级数
\[
\sum_{r\ge 1}\log\frac{1+a_r}{1-a_r}
\]
收敛，故 \(\mathcal K_\varphi\) 有限。

最后，由定义
\[
(\mathcal A_R\delta)_m=\sum_{j=0}^{R}\beta_j^{(R)}\delta_{m+j},
\]
故
\[
\bigl|(\mathcal A_R\delta)_m\bigr|
\le
\sum_{j=0}^{R}\bigl|\beta_j^{(R)}\bigr|\,\max_{0\le j\le R}|\delta_{m+j}|.
\]
将上式代入有限阶闭式并利用其对 \(R\) 单调递增且上界为 \(\mathcal K_\varphi\)，便得到所示扰动界。证毕。
\end{proof}

\begin{theorem}[二次 \texorpdfstring{$f$}{f}-散度标量的自校准最小充分窗]\label{thm:conclusion-foldgauge-pi-selfcalibrated-minimal-window}
取
\[
f_2(t):=(t-1)^2,
\qquad
D_2^\infty:=D_{f_2}^\infty.
\]
则有显式反演
\[
D_2^\infty=\frac{\varphi^3+1}{5}-1=\frac{2\varphi-3}{5},
\qquad
\varphi=\frac{5D_2^\infty+3}{2},
\qquad
q=\frac{5D_2^\infty+3}{4}.
\]
从而对任意 \(R\ge 1\)，加速多项式 \(P_R\) 都可写成 \(D_2^\infty\) 的有理--代数组合：
\[
P_R(z)
=
\prod_{r=1}^{R}
\frac{
z-\left(\frac{5D_2^\infty+3}{4}\right)^{2r-1}
}{
1-\left(\frac{5D_2^\infty+3}{4}\right)^{2r-1}
}.
\]
因此，在固定窗仿射对数类 \(\mathfrak A_R\) 中，数据组
\[
\mathcal S_{m,R}
:=
\bigl(D_2^\infty;\ C_m,C_{m+1},\dots,C_{m+R}\bigr)
\]
构成达到误差阶 \(q^{(2R+1)m}\) 的 \(\pi\) 内生恢复所需的最小充分窗：它足以唯一给出
\[
\Pi_m^{\langle R\rangle},
\]
且不能再缩短为仅含 \(R\) 个连续 \(C\)-值而仍保持同阶最优性。
\end{theorem}
\begin{proof}
结论 \ref{con:xi-time-part9l-chi2-generates-power-divergence-and-cumulants} 已给出
\[
D_2^\infty=\frac{2\varphi-3}{5},
\]
故 \(\varphi\) 与 \(q=\varphi/2\) 都由单一标量 \(D_2^\infty\) 唯一确定。将其代入定理 \ref{thm:conclusion-foldgauge-pi-oddlayer-annihilation-hierarchy} 中 \(P_R\) 的定义，即得所示显式公式。

于是，给定 \(\mathcal S_{m,R}\) 后，\(q\)、多项式 \(P_R\)、算子 \(\mathcal A_R\) 与估计器 \(\Pi_m^{\langle R\rangle}\) 均被唯一决定，从而 \(\mathcal S_{m,R}\) 在固定窗仿射对数类中是充分的。

最小性来自定理 \ref{thm:conclusion-foldgauge-pi-affine-log-accelerator-uniqueness}：若只保留 \(R\) 个连续 \(C\)-值，则相应仿射对数窗口长度至多为 \(R\)，不可能达到 \(O(q^{(2R+1)m})\) 的误差阶。证毕。
\end{proof}

\begin{theorem}[偶值 \texorpdfstring{$\zeta$}{zeta} 的有限窗抽取公式]\label{thm:conclusion-foldgauge-even-zeta-finite-window-extraction}
\leanverified{paper\_conclusion\_foldgauge\_even\_zeta\_finite\_window\_extraction}
对每个整数 \(r\ge 1\)，约定 \(P_0\equiv 1\)，并定义
\[
\mathcal R_{m,r}
:=
q^{-(2r-1)m}
\bigl(P_{r-1}(T)L\bigr)_m
=
q^{-(2r-1)m}
\bigl(P_{r-1}(T)(\log C-\log(2\pi))\bigr)_m.
\]
则极限存在，且满足
\[
\mathcal R_{m,r}
=
a_r\,P_{r-1}(q^{2r-1})
+
O(q^{2m}).
\]
从而有显式抽取公式
\[
\zeta(2r)
=
(-1)^{r-1}
\frac{(2\pi)^{2r}}{2(2r)!}\,
\frac{r(2r-1)}{\varphi^{-1}+\varphi^{2r-3}}\,
\frac{1}{P_{r-1}(q^{2r-1})}\,
\lim_{m\to\infty}\mathcal R_{m,r}.
\]
因此，第 \(r\) 个偶值 \(\zeta(2r)\) 可由长度恰为 \(r\) 的连续尺度窗
\[
(C_m,C_{m+1},\dots,C_{m+r-1})
\]
直接抽取；其误差统一为 \(O(q^{2m})\)。又由于 \(2\pi=\lim_{m\to\infty}C_m\) 已由同一通道内生恢复，故右端不需任何外部算术输入。
\end{theorem}
\begin{proof}
把定理 \ref{thm:conclusion-foldgauge-pi-oddlayer-annihilation-hierarchy} 中的 \(R\) 替换为 \(r-1\)，并注意该定理作用于 \(L_m=\log C_m-\log(2\pi)\) 时常数项已经消失，得到
\[
\bigl(P_{r-1}(T)L\bigr)_m
=
a_r P_{r-1}(q^{2r-1})q^{(2r-1)m}
+
O\!\bigl(q^{(2r+1)m}\bigr).
\]
两边同乘 \(q^{-(2r-1)m}\) 便得
\[
\mathcal R_{m,r}
=
a_r P_{r-1}(q^{2r-1})+O(q^{2m}).
\]
由 \(q^{2r-1}\) 不属于 \(P_{r-1}\) 的零点集，可知
\[
P_{r-1}(q^{2r-1})\neq 0,
\]
故极限存在且可反解出 \(a_r\)。

最后，代入
\[
a_r=
\frac{B_{2r}}{r(2r-1)}(\varphi^{-1}+\varphi^{2r-3}),
\qquad
B_{2r}=(-1)^{r-1}\frac{2(2r)!}{(2\pi)^{2r}}\zeta(2r),
\]
整理即得所示显式抽取公式。证毕。
\end{proof}

\begin{theorem}[折叠误差正规形的 Koenigs 型与 AGM 非共轭性]\label{thm:conclusion-foldgauge-pi-koenigs-vs-agm-nonconjugacy}
记
\[
e_m:=\log C_m-\log(2\pi),
\qquad
q:=\frac{\varphi}{2}.
\]
则 fold-gauge \(\pi\) 误差满足局部正规形
\[
e_{m+1}
=
q\,e_m+\Gamma e_m^3+O(e_m^5),
\]
其中
\[
\Gamma
=
\frac{q^3-q}{a_1^3}\,a_2
=
\frac{3(2+\sqrt5)}{32}\neq 0,
\qquad
a_1=\frac{1}{3\varphi},
\qquad
a_2=\frac{B_4}{2\cdot 3}(\varphi^{-1}+\varphi)
=
-\frac{\sqrt5}{180}.
\]
因此折叠 \(\pi\) 通道的局部动力学属于非零乘子的 Koenigs 型机制。另一方面，经典 AGM/Gauss--Salamin 误差迭代具有超吸引正规形
\[
u_{n+1}=\kappa u_n^2+O(u_n^3),
\qquad
\kappa\neq 0,
\]
其乘子为 \(0\)，属于 B\"ottcher 型机制。由于局部形式共轭与局部解析共轭都保持线性项乘子，二者不可能共轭。因而 fold-gauge 通道的自然加速并不是 AGM 式平方化，而只能沿定理 \ref{thm:conclusion-foldgauge-pi-oddlayer-annihilation-hierarchy} 的奇层消元塔发生。
\end{theorem}
\begin{proof}
由定理 \ref{thm:xi-time-part64a-fold-gauge-quintic-renormalization-normal-form}，
\[
e_{m+1}
=
q e_m+\frac{3(2+\sqrt5)}{32}e_m^3+O(e_m^5).
\]
另一方面，定理 \ref{thm:fold-bin-gauge-constant-stirling-bernoulli-hierarchy} 在前两层给出
\[
e_m=a_1q^m+a_2q^{3m}+O(q^{5m}),
\]
其中 \(a_1=1/(3\varphi)\)、\(a_2=-\sqrt5/180\)。比较三次项便得到
\[
\Gamma=\frac{q^3-q}{a_1^3}a_2=\frac{3(2+\sqrt5)}{32}\neq 0.
\]
再由定理 \ref{thm:xi-time-part69-fold-gauge-formal-koenigs-linearization}，该重整化确属 Koenigs 线性化机制。

对 AGM/Gauss--Salamin 迭代，经典误差理论给出其局部正规形为
\[
u\longmapsto \kappa u^2+O(u^3),
\]
故其线性项乘子为 \(0\)。若存在局部形式共轭或局部解析共轭 \(h\) 把一个映射 \(F\) 共轭到另一个映射 \(G\)，则
\[
h\circ F=G\circ h,
\qquad
h(0)=0,
\qquad
h'(0)\neq 0.
\]
在 \(0\) 点求导得到
\[
h'(0)F'(0)=G'(0)h'(0),
\]
故 \(F'(0)=G'(0)\)。因此乘子 \(q\neq 0\) 的 fold-gauge 正规形不可能与乘子 \(0\) 的 AGM 正规形共轭。证毕。
\end{proof}

\noindent
因此，fold-gauge 标量 \(\pi\) 通道所承载的并不是一条单纯线性收敛的数值序列，而是一条严格可审计的奇次 Bernoulli 谱链：它既能被唯一最优且一致有界条件的有限窗算子逐层湮灭，又能由单一 \(\chi^2\) 极限常数完成全部校准，并进一步把偶值 \(\zeta\)-塔压缩为显式有限窗读出；与此同时，其局部动力学在正规形层面与 AGM 的超吸引平方化机制严格分离。于是，折叠 \(\pi\) 估计的自然加速原理不再是外加模性输入或额外位地址，而是项目内部本已固有的奇层消元、谱抽取与 Koenigs 线性化三条结构同时闭合。

\endinput
